% Pursuing Stacks (cotes 134-2 to 134-8): archive cote 134-6, site page 18 (PDF page 19), typescript page 337.
% typescript. Draft compared: ps5.tex:2279-2334.
% AI-transcribed by Claude Opus 5.5 (claude-opus-5-5, adaptive thinking, effort high), Message Batches API, batch pilot (msgbatch_0189Nhf1xB615YWnztubtLfc), 2026-11-12; reviewed against both supplied scan halves; see CHECK_DECISIONS.tsv.
\scanpage{134-6/18}
\mspage{337}

\deleted{properties} are stable under taking \xout{pro binary} products (of finite non-%
empty families of categories $A_i$), hence in terms of the three standard
\over[typed]{t}{c}ases, the possibility of satisfying them too by the ``multistandard''
test categories. I wonder if there are any other ways (up to equiva-%
lence). If we take categories such as $\hand{\widetilde{\Delta}}$ (non ordered simplices) we
still get contractors, but objects have non trivial automo\xout{t}rphisms, and
if we take categories such as $\hand{\Delta}^{f}$ (ordered simplices without degenera-%
cy operations, only boundary maps) it is true that isomorphisms are
identities, but the category is no longer a test category \xout{not even}\inserted[typed]{but only} a
weak one.

If we do not insist on the rigidity assumption (isomorphisms are
identities), but on suitable choice of so called ``\uline{strict} boundary
operations'' within A , then it would seem after all that we do have
a lot more elbow freedom then it seemed, by the end of our reflections
on that matter two days ago (cf. p. 318), where the picture of the re-%
levant data and corresponding \xout{structur} construction of chain complexes
was still a little confused. Let now A be the category called $A_o$ ih
loc. cit. We don't have to modify it in order to introduce orientati-%
ons of cells i(a) as extra structure and take account of this in defi-%
ning \xout{new} a new notion of map\hand{s.} Therefore it is clear that A \uline{just as it}
\uline{i\over[typed]{t}{s}, is a strict test category} (\xout{but} presumably \uline{not} a contractor though).
T$\overset{\illegible{1}}{h}$ere is the problem of course of isomorphisms which are not identi-%
ties, and particularly of non trivial automorphisms - for instance
the object I (playing the part of the unit segment) has a non trivial
automorphism, the elimination of which does not look so trivial !
However, there \uline{is} a rather evident way of cutting out \uline{strict} bounda-%
ry maps, in a way as to satisfy the transitivity condition of p.335
- namely by taking boundary maps (13) $\hand{\partial}: b \hand{\to} a$ which a\over[typed]{r}{r}e \uline{inclusions}
in the strict sense, namely the inclusion map of a\xout{n} \uline{subset} of a , en-%
dowed with the induced order relation.

Thus there are many other \xout{ways}\inserted[typed]{cases} still than just multi-standard
test categories for getting a canonical chain complex $L.$ in Add($A^o$),
\deleted{so that} \inserted{and for whom now} the question makes sense as to whether $L.$ is an abelianiz\over{\unsure{o|e}r}{ator}.
In the construction above, we were careful to assume that the full
subcategory \xout{a} A of (Ord), besides containing I , was stable under
finite products, so as to make sure it comes out as a test category.
T$\overset{\illegible{1}}{\underset{\illegible{1}}{h}}$e silly thing is that this condition is not satisfied by any one among
the standard test categories - thus it seems reasonable to try and

\margin{18}

% Visual collation: both supplied image halves inspected; see CHECK_DECISIONS_134_7_8.tsv.
